% requires \usepackage{tabularx} and \usepackage{booktabs} in the preamble
\begin{table}[t]
\centering
{\footnotesize
\begin{tabularx}{\linewidth}{@{} >{\raggedright\arraybackslash}X @{}}
\toprule
\textbf{Original sentence.} Substituting the reparameterization (...) into the preference model, the partition function cancels, and we can express the human preference probability in terms of only the optimal policy $\pi^*$ and reference policy $\pi_\text{ref}$. \tabularnewline
\midrule
\textbf{Blog.} Direct Preference Optimization (DPO) offers a neat change of perspective. Instead of treating reward modeling and policy optimization as separate steps, DPO re-parameterizes the reward class so the optimal policy corresponding to a reward has a closed form. That lets you fit the policy directly to human comparisons with a simple classification loss --- no RL required. \tabularnewline
\midrule
\textbf{Stack Exchange.} \textbf{Q:} How exactly does the partition function cancel out in the DPO derivation? They say the partition function term $\beta \log Z(x)$ cancels --- can someone show the algebraic steps explicitly? \textbf{A:} Short answer: because the partition function term in the reparameterization depends only on the prompt $x$, it is identical for both completions and therefore subtracts out when you form the reward difference that appears inside the Bradley--Terry sigmoid. \tabularnewline
\midrule
\textbf{Textbook.} The partition function term $\beta\log Z(x)$ depends only on $x$, so it cancels when forming differences. Preference models like Bradley--Terry depend only on reward differences, not on absolute reward values. Therefore the unknown normalization that made reward recovery hard is irrelevant to the likelihood of observed pairwise comparisons. This is precisely why we can reparameterize rewards in terms of policies and then optimize policy parameters directly with a standard cross-entropy loss. \tabularnewline
\bottomrule
\end{tabularx}
}
\caption{\small{To study how the \textit{format} of supporting text affects knowledge acquisition, we rewrite each source paper into three explanatory genres while preserving the underlying facts. Shown here for a single sentence from `Direct Preference Optimization: Your Language Model is Secretly a Reward Model', the \textit{blog} adopts an accessible, narrative tone; the \textit{Stack Exchange} entry frames the knowledge as a question-and-answer exchange; and the \textit{textbook} situates it in formal, structured exposition. Each genre develops the same machinery from a different angle: the blog pitches DPO's reparameterization of reward as policy, the Stack Exchange entry works through why the partition function cancels, and the textbook explains why that function is intractable to compute directly.}}
\label{tab:explanation_formats}
\vspace{-5mm}
\end{table}
